% GENERATED FILE. Edit canonical lessons, then run npm run generate:books.
% canonical-source-hash: fnv1a64:16fb4e3e685f3355
% canonical-lessons: TE-C515-medadu, TE-C515-naram, TE-C515-kaleyam, TE-C515-vennemuka, TE-C515-svasa

\chapter{Inside the Body}
\label{ch:te-inside-the-body}
\hlchaptermodality{515}

\begin{chapteropening}
\textbf{By the end of this chapter:} \emph{I can name the brain, nerves, liver, spine and breath.}

Complete the last lesson of chapter 515: I can name the brain, nerves, liver, spine and breath.
\end{chapteropening}

\section[\texorpdfstring{\te{మెదడు} (medaḍu)}{మెదడు (medaḍu)}]{\te{మెదడు} (medaḍu) — the brain}

\label{lesson:TE-C515-medadu}



\begin{quote}
Before the new one: say the Telugu for to groan, then the Telugu for to faint.
\end{quote}

\subsection*{You'll want to know: \te{మెదడు}}
\textbf{\te{మెదడు}} — \emph{medaḍu} — \textquotedblleft{}the brain\textquotedblright{}.

\begin{quote}
\te{మెదడు} \te{శరీరాన్ని} \te{నడిపిస్తుంది}.
\end{quote}

\begin{quote}
\emph{medaḍu śarīrānni naḍipistundi.}
\end{quote}

\begin{quote}
The brain runs the body.
\end{quote}

Also your wits: \te{మెదడు} \te{వాడు}, to use your head.

\begin{cousinweb}
\textbf{\te{మెదడు}} (\emph{medaḍu}) is native Dravidian: DEDR 5062 lists Telugu \emph{medaḍu}, \textquotedblleft{}brain\textquotedblright{}, with Kannada \emph{meduḷ} and Kota \emph{medl}.
\end{cousinweb}

\begin{grammarlens}[title={What you've built}]
One more word for describing what you lived and felt.
\end{grammarlens}

\subsection*{Your turn}
\begin{itemize}
\raggedright
  \item \emph{Say it:} \emph{medaḍu}
  \item \emph{Say it:} \emph{medaḍu}, once more
  \item \emph{Say it:} \emph{medaḍu śarīrānni naḍipistundi.} — The brain runs the body.
  \item \emph{Recall:} say the Telugu for to groan, then the Telugu for to faint, then say \emph{medaḍu} again
\end{itemize}

\begin{tcolorbox}[breakable,colback=teal!4,colframe=teal!35!black,title={Before you move on}]
What does \te{మెదడు} mean? (\textquotedblleft{}The brain\textquotedblright{}.) Say it once more.
\end{tcolorbox}
\section[\texorpdfstring{\te{నరం} (naraṁ)}{నరం (naraṁ)}]{\te{నరం} (naraṁ) — a nerve; a vein}

\label{lesson:TE-C515-naram}



\begin{quote}
Before the new one: say the Telugu for to faint, then the Telugu for the brain.
\end{quote}

\subsection*{You'll want to know: \te{నరం}}
\textbf{\te{నరం}} — \emph{naraṁ} — \textquotedblleft{}a nerve; a vein\textquotedblright{}.

\begin{quote}
\te{నర్సు} \te{నా} \te{చేతి} \te{నరంలో} \te{సూది} \te{వేసింది}.
\end{quote}

\begin{quote}
\emph{narsu nā cēti naraṁlō sūdi vēsindi.}
\end{quote}

\begin{quote}
The nurse put the needle into a vein in my arm.
\end{quote}

Plural: \te{నరాలు}, \emph{narālu}.

Plural: \te{నరాలు}.

\begin{cousinweb}
\textbf{\te{నరం}} (\emph{naraṁ}) is native Dravidian: DEDR 2903 lists Telugu \emph{naramu}, \textquotedblleft{}vein, artery, nerve, tendon\textquotedblright{}, with Tamil \emph{narampu}.
\end{cousinweb}

\begin{grammarlens}[title={What you've built}]
One more word for describing what you lived and felt.
\end{grammarlens}

\subsection*{Your turn}
\begin{itemize}
\raggedright
  \item \emph{Say it:} \emph{naraṁ}
  \item \emph{Say it:} \emph{naraṁ}, once more
  \item \emph{Say it:} \emph{narsu nā cēti naraṁlō sūdi vēsindi.} — The nurse put the needle into a vein in my arm.
  \item \emph{Recall:} say the Telugu for to faint, then the Telugu for the brain, then say \emph{naraṁ} again
\end{itemize}

\begin{tcolorbox}[breakable,colback=teal!4,colframe=teal!35!black,title={Before you move on}]
What does \te{నరం} mean? (\textquotedblleft{}A nerve; a vein\textquotedblright{}.) Say it once more.
\end{tcolorbox}
\section[\texorpdfstring{\te{కాలేయం} (kālēyaṁ)}{కాలేయం (kālēyaṁ)}]{\te{కాలేయం} (kālēyaṁ) — the liver}

\label{lesson:TE-C515-kaleyam}



\begin{quote}
Before the new one: say the Telugu for the brain, then the Telugu for a nerve.
\end{quote}

\subsection*{You'll want to know: \te{కాలేయం}}
\textbf{\te{కాలేయం}} — \emph{kālēyaṁ} — \textquotedblleft{}the liver\textquotedblright{}.

\begin{quote}
\te{డాక్టరు} \te{నా} \te{కాలేయాన్ని} \te{పరీక్షించారు}.
\end{quote}

\begin{quote}
\emph{ḍākṭaru nā kālēyānni parīkṣiñcāru.}
\end{quote}

\begin{quote}
The doctor examined my liver.
\end{quote}

A -\te{ం} noun takes -\te{ా}- before an ending: \te{కాలేయానికి}, \te{కాలేయాన్ని}.

\begin{cousinweb}
\textbf{\te{కాలేయం}} (\emph{kālēyaṁ}) is borrowed from Sanskrit \emph{kāleya}, \textquotedblleft{}the liver\textquotedblright{} (Monier-Williams).
\end{cousinweb}

\begin{grammarlens}[title={What you've built}]
One more word for describing what you lived and felt.
\end{grammarlens}

\subsection*{Your turn}
\begin{itemize}
\raggedright
  \item \emph{Say it:} \emph{kālēyaṁ}
  \item \emph{Say it:} \emph{kālēyaṁ}, once more
  \item \emph{Say it:} \emph{ḍākṭaru nā kālēyānni parīkṣiñcāru.} — The doctor examined my liver.
  \item \emph{Recall:} say the Telugu for the brain, then the Telugu for a nerve, then say \emph{kālēyaṁ} again
\end{itemize}

\begin{tcolorbox}[breakable,colback=teal!4,colframe=teal!35!black,title={Before you move on}]
What does \te{కాలేయం} mean? (\textquotedblleft{}The liver\textquotedblright{}.) Say it once more.
\end{tcolorbox}
\section[\texorpdfstring{\te{వెన్నెముక} (vennemuka)}{వెన్నెముక (vennemuka)}]{\te{వెన్నెముక} (vennemuka) — the spine; the backbone}

\label{lesson:TE-C515-vennemuka}



\begin{quote}
Before the new one: say the Telugu for a nerve, then the Telugu for the liver.
\end{quote}

\subsection*{You'll want to know: \te{వెన్నెముక}}
\textbf{\te{వెన్నెముక}} — \emph{vennemuka} — \textquotedblleft{}the spine; the backbone\textquotedblright{}.

\begin{quote}
\te{పడిపోవడంతో} \te{తాతయ్య} \te{వెన్నెముకకు} \te{దెబ్బ} \te{తగిలింది}.
\end{quote}

\begin{quote}
\emph{paḍipōvaḍantō tātayya vennemukaku debba tagilindi.}
\end{quote}

\begin{quote}
Grandpa hurt his spine in a fall.
\end{quote}

\te{వెన్ను}, back, + \te{ఎముక}, bone. As in English, it also means the mainstay.

\begin{cousinweb}
\textbf{\te{వెన్నెముక}} (\emph{vennemuka}) is native Dravidian: \emph{vennu}, \textquotedblleft{}the back\textquotedblright{} (DEDR 5488), and \emph{emuka}, \textquotedblleft{}bone\textquotedblright{}, which DEDR 839 lists with Tamil \emph{elumpu}.
\end{cousinweb}

\begin{grammarlens}[title={What you've built}]
One more word for describing what you lived and felt.
\end{grammarlens}

\subsection*{Your turn}
\begin{itemize}
\raggedright
  \item \emph{Say it:} \emph{vennemuka}
  \item \emph{Say it:} \emph{vennemuka}, once more
  \item \emph{Say it:} \emph{paḍipōvaḍantō tātayya vennemukaku debba tagilindi.} — Grandpa hurt his spine in a fall.
  \item \emph{Recall:} say the Telugu for a nerve, then the Telugu for the liver, then say \emph{vennemuka} again
\end{itemize}

\begin{tcolorbox}[breakable,colback=teal!4,colframe=teal!35!black,title={Before you move on}]
What does \te{వెన్నెముక} mean? (\textquotedblleft{}The spine; the backbone\textquotedblright{}.) Say it once more.
\end{tcolorbox}
\section[\texorpdfstring{\te{శ్వాస} (śvāsa)}{శ్వాస (śvāsa)}]{\te{శ్వాస} (śvāsa) — breath; breathing}

\label{lesson:TE-C515-svasa}



\begin{quote}
Before the new one: say the Telugu for the liver, then the Telugu for the spine.
\end{quote}

\subsection*{You'll want to know: \te{శ్వాస}}
\textbf{\te{శ్వాస}} — \emph{śvāsa} — \textquotedblleft{}breath; breathing\textquotedblright{}.

\begin{quote}
\te{లోతుగా} \te{శ్వాస} \te{తీసుకుంటే} \te{మనసు} \te{ప్రశాంతంగా} \te{ఉంటుంది}.
\end{quote}

\begin{quote}
\emph{lōtugā śvāsa tīsukuṇṭē manasu praśāntaṅgā uṇṭundi.}
\end{quote}

\begin{quote}
If you breathe deeply, the mind is calm.
\end{quote}

\te{శ్వాస} \te{తీసుకో}, to take a breath.

\begin{cousinweb}
\textbf{\te{శ్వాస}} (\emph{śvāsa}) is borrowed from Sanskrit \emph{śvāsa}, \textquotedblleft{}breath, breathing\textquotedblright{} (CDIAL 12769).
\end{cousinweb}

\begin{grammarlens}[title={What you've built}]
One more word for describing what you lived and felt.
\end{grammarlens}

\subsection*{Your turn}
\begin{itemize}
\raggedright
  \item \emph{Say it:} \emph{śvāsa}
  \item \emph{Say it:} \emph{śvāsa}, once more
  \item \emph{Say it:} \emph{lōtugā śvāsa tīsukuṇṭē manasu praśāntaṅgā uṇṭundi.} — If you breathe deeply, the mind is calm.
  \item \emph{Recall:} say the Telugu for the liver, then the Telugu for the spine, then say \emph{śvāsa} again
\end{itemize}

\begin{tcolorbox}[breakable,colback=teal!4,colframe=teal!35!black,title={Before you move on}]
What does \te{శ్వాస} mean? (\textquotedblleft{}Breath; breathing\textquotedblright{}.) Say it once more.
\end{tcolorbox}
